\subsection{Behold the Handmaid: the Ember Days}\label{sec:ember}

On the Wednesday after the Third Sunday the Church reads the Gospel of the
Annunciation. The station is St Mary Major. Isaiah has prepared the way in the
lessons before it: ``Behold a virgin shall conceive, and bear a son and his name
shall be called Emmanuel.'' Then the angel Gabriel is sent to Nazareth, greets
Mary, and tells her what God will do.

St Bernard preached a series of homilies on this Gospel, and in the fourth he
stands beside the angel and waits with him for her reply. He brings before Mary
the whole of mankind: Adam weeping in exile from Paradise, Abraham and David,
the holy fathers still dwelling in the region of the shades of death, the whole
race prostrate at her feet. Then he speaks for them himself:

\begin{quote}
For lo! the price of our salvation is now offered to thee: if thou wilt only
consent, we shall at once be set at liberty. We have been created by the
eternal Word of God, and behold we die: by thy momentary word we must be renewed
and restored to life. \dots\ Make haste, therefore, to answer the Angel, or
rather to answer the Lord through the Angel. Say the word and receive the Word.
Utter thy human word and conceive the Divine Word. Pronounce the transitory word
and embrace the Word everlasting.\par
{\smallskip\noindent\small ---St Bernard, Fourth Sermon on the Virgin Mother\par}
\end{quote}

And she says it:

\begin{quote}
And Mary said: Behold the handmaid of the Lord: be it done to me according to
thy word. And the angel departed from her.\par
{\smallskip\noindent\small ---Luke 1:38\par}
\end{quote}

The Word who made the world waited for the word of a virgin of Nazareth, and
with her \latin{Fiat} the first coming began. Those who say the Angelus repeat
her words morning, noon and evening; on this Wednesday the Church sets before
them the whole Gospel the words come from.

\subsubsection*{The December fast}

The Ember Days are the Roman fast of the tenth month on which St Leo the Great
preached, still kept in December on Wednesday, Friday and Saturday. Leo gave
his people a reason for fasting at this time of year that is easy to forget in
a city:

\begin{quote}
Dearly beloved brethren, with the care which becometh us as the shepherd of
your souls, we urge upon you the rigid observance of this December Fast. The
month of December hath come round again, and with it this devout custom of the
Church. The fruits of the year, which is drawing to a close, are now all
gathered in, and we most meetly offer our abstinence to God as a sacrifice of
thanksgiving. And what can be more useful than fasting, that exercise by which
we draw nigh to God, make a stand against the devil, and overcome the softer
enticements of sin?\par
{\smallskip\noindent\small ---St Leo the Great, Sermon 13\par}
\end{quote}

The harvest is in, and the Church gives thanks for it by going without. Keep
the three days, as far as you can, with the fast, the prayer and the alms Leo
asked for, and go to Mass on one of them if there is a Mass near you.

\subsubsection*{Friday: the leap in the womb}

On Friday the station is the church of the Twelve Apostles. Isaiah's root of
Jesse is read, and then the Gospel of the Visitation. Mary goes into the hill
country to her cousin Elizabeth, and when Elizabeth hears her greeting ``the
infant leaped in her womb.'' The unborn John, who will be the voice in the
wilderness, makes his first sign before the unborn Word. St Ambrose, in the
lesson read at Matins, looks at the road Mary took:

\begin{quote}
She went into the hill country in joyful haste. And is it not something that
she went up into the hills? God was already in her womb, and her feeling bore
her continually upward. The grace of the Holy Spirit knoweth no slow
working.\par
{\smallskip\noindent\small ---St Ambrose, On St Luke, Book 2\par}
\end{quote}

On the day the Gospel shows Christ carried hidden in his mother over the hills,
the Church takes up the First Sunday's words again and asks him to come:

\begin{studyframe}
\latin{Excita, quaesumus, Domine, potentiam tuam, et veni: ut hi, qui in tua
pietate confidunt, ab omni citius adversitate liberentur: Qui vivis et
regnas cum Deo Patre in unitate.}
\medskip

Bestir, O Lord, Thy might, we pray Thee, and come; that these who trust in Thy
kindness and love, may the more speedily be delivered from all adversity. Who
livest.\par
\end{studyframe}

\subsubsection*{Saturday: dew in the furnace}

On Saturday the station is St Peter's, where Leo bade his people keep vigil
with the Apostle. The Mass of the day has a shorter form and a longer, and the
longer is for the conventual Mass and for the Mass at which Holy Orders are
conferred: the Ember Saturday is a day of ordinations. In
the longer form the Church reads four lessons from Isaiah, a lesson from Daniel,
St Paul to the Thessalonians, and the Gospel of John's preaching in the
fifteenth year of Tiberius. The Daniel lesson tells of the young men in the
furnace:

\begin{quote}
But the angel of the Lord went down with Azarias and his companions into the
furnace: and he drove the flame of the fire out of the furnace, and made the
midst of the furnace like the blowing of a wind bringing dew, and the fire
touched them not at all, nor troubled them, nor did them any harm.\par
{\smallskip\noindent\small ---Daniel 3:49--50\par}
\end{quote}

On the eve of a Sunday whose Mass begins ``Drop down dew, ye heavens'', the
Church reads of a dew that came down into the furnace and kept Azarias and his
companions alive.

Pray this Saturday for those who receive Holy Orders, and for the priests who
will hear your confession and say the Masses of Christmas.
